\chapter{Complexity — Unary Code}
%%% generated by the blueprint pipeline from TCSlib.Complexity.CircuitComplexity.UnaryCode
\section{Overview}

Declarations of the Lean module \texttt{TCSlib.Complexity.CircuitComplexity.UnaryCode}.

\section{Declarations}

\begin{definition}[Unary encoding of a natural number]
\lean{BoolCircuit.encodeNat}
\label{BoolCircuit.encodeNat}
\leanok
The unary encoding of a natural number $k$ is the bit string consisting of $k$ copies of
$1$, followed by a single $0$.
\end{definition}

\begin{definition}[Decoder for natural number strings]
\lean{BoolCircuit.decodeNat}
\label{BoolCircuit.decodeNat}
\leanok
A decoder on finite Boolean strings returns no value for the empty string. For a
nonempty string, it counts the initial consecutive $1$ bits up to the first $0$ and
returns that count together with the remaining suffix; if no $0$ occurs, it returns no
value.
\end{definition}

\begin{theorem}[Natural number encoding and decoding]
\lean{BoolCircuit.decodeNat_encodeNat}
\label{BoolCircuit.decodeNat_encodeNat}
\leanok
\uses{BoolCircuit.decodeNat, BoolCircuit.encodeNat}
\difficulty{4}
For every natural number $k$ and Boolean-list suffix $r$, decoding the concatenation of
$k$ true bits, one false bit, and $r$ returns $(k,r)$.
\end{theorem}
